\chapter{\texttt{detector.common} --- generic event-rate folding infrastructure}
\label{ch:detector-common}

Every propagation chapter so far ends at a \emph{true} quantity: an
oscillated flavour flux, or a survival probability $P_{ee}(E_\nu)$, defined
at the level of the neutrino itself. No real detector measures that
directly. What is actually recorded is a finite number of counts in bins of
a \emph{reconstructed} observable (an electron recoil energy, a photon
count, a charge deposit), after the true flux has been filtered through a
target's interaction cross section, the detector's finite energy
resolution, its selection efficiency, and its finite exposure. This chapter
covers the module that performs exactly that filtering, generically, for
any detector: \texttt{detector.common}.

Mirroring \texttt{core.common.hamiltonian}'s ``single Hamiltonian assembly
module'' principle (\cref{ch:core-common}), \texttt{detector.common.event\_rate}
is the \emph{only} event-rate folding code in the project: every concrete
detector's own \texttt{event\_rate.py} (currently only
\texttt{detector.borexino}, \cref{ch:detector-borexino}) composes the same
four functions defined here with its own target, response, efficiency, and
exposure, rather than re-deriving the folding formula. The interaction
cross sections folded in are supplied by \texttt{detector.interaction}
(\cref{ch:detector-interaction}); this chapter treats them as an external
input, $d\sigma/dT(E_\nu,T)$, of unspecified physical origin.

\section{\texttt{event\_rate.py} --- flux to predicted counts}
\label{sec:detector-event-rate}

\modulehead{detector/common/event\_rate.py}{the four-step forward-folding
chain: cross-section convolution, energy-response smearing, efficiency,
and exposure-weighted binning.}

\begin{theorybox}
\subsection*{Step 1: true event rate from flux and cross section}
A neutrino flux $\Phi(E_\nu)$ (differential in true neutrino energy,
$\mathrm{cm^{-2}\,s^{-1}\,MeV^{-1}}$) incident on $N_{\rm target}$ target
particles produces a differential interaction rate in the true
observable $T$ (e.g. an electron recoil kinetic energy) by the standard
flux $\times$ cross section $\times$ target-count folding, integrated over
the incident neutrino energy that can produce a given $T$. Only the $\nu_e$
versus $(\nu_\mu+\nu_\tau)$ split of the flux matters, since every
interaction currently in \texttt{detector.interaction} treats $\nu_\mu$ and
$\nu_\tau$ identically (\cref{ch:detector-interaction}), so a full
three-flavour flux collapses to two effective channels $\Phi_e,\Phi_x$:
\begin{equation}
  \keyresult{\frac{dR}{dT}(T) = N_{\rm target}\int dE_\nu\,
    \left[\Phi_e(E_\nu)\,\frac{d\sigma_e}{dT}(E_\nu,T)
        + \Phi_x(E_\nu)\,\frac{d\sigma_x}{dT}(E_\nu,T)\right].}
  \label{eq:detector-true-rate}
\end{equation}
The integral runs over the full $E_\nu$ grid supplied; the cross sections
themselves enforce the kinematic bound $T\leq T_{\max}(E_\nu)$ by vanishing
outside it (\cref{ch:detector-interaction}), so no explicit integration
limit needs to be applied here.

\subsection*{Step 2: finite energy resolution (response/migration)}
No real detector measures $T$ exactly: photostatistics, calibration
uncertainty, and reconstruction algorithms smear the true deposited energy
$T$ into a reconstructed observable $T'$. This is encoded by a
\emph{response} (or migration) function $R(T'|T)$ -- the conditional
probability density of reconstructing $T'$ given a true deposit $T$,
normalised as $\int dT'\,R(T'|T)=1$ for every $T$ (probability
conservation: every true event is reconstructed \emph{somewhere}). The
reconstructed-energy spectrum follows by convolution,
\begin{equation}
  \keyresult{\frac{dR}{dT'}(T') = \int dT\,R(T'|T)\,\frac{dR}{dT}(T).}
  \label{eq:detector-response}
\end{equation}
The concrete functional form of $R(T'|T)$ is detector-specific (a Gaussian
for a photoelectron-counting scintillator, \cref{sec:detector-response}
below; a Monte-Carlo migration matrix for a Cherenkov detector) and is not
fixed by this module.

\subsection*{Step 3: detection/selection efficiency}
Analysis cuts (fiducial-volume selection, trigger threshold,
pulse-shape discrimination, ...) act on what is actually measured, i.e.\ on
the reconstructed observable $T'$, not on the inaccessible true $T$. The
efficiency $\varepsilon(T')\in[0,1]$ is therefore applied \emph{after} the
response convolution, not before:
\begin{equation}
  \keyresult{\frac{dR_{\rm det}}{dT'}(T') = \varepsilon(T')\,\frac{dR}{dT'}(T').}
  \label{eq:detector-efficiency}
\end{equation}

\subsection*{Step 4: exposure-weighted binning, plus background}
Finally, the detected differential spectrum is integrated over each finite
analysis bin $[T'_i,T'_{i+1}]$ of the published/observed spectrum and scaled
by the exposure (live time, optionally combined with a normalisation
convention matching a specific published rate table) to give an absolute
predicted count $N_i$ -- the quantity actually compared against data in a
Poisson likelihood:
\begin{equation}
  \keyresult{N_i = {\rm exposure}\int_{T'_i}^{T'_{i+1}} dT'\,
    \frac{dR_{\rm det}}{dT'}(T') \;+\; N_i^{\rm bkg}.}
  \label{eq:detector-bin-counts}
\end{equation}
The optional background term $N_i^{\rm bkg}$ is added, not folded into the
convolution: it is modelled as an independent, incoherent process (radioactive
contaminants, cosmogenic isotopes, external gammas, ...) uncorrelated with
the oscillated-flux signal being propagated through
\labelcref{eq:detector-true-rate,eq:detector-response,eq:detector-efficiency},
consistent with the standard ``signal + background'' template used in a
counting experiment's likelihood. See \texttt{background.py} below for how
this term is (not yet) populated.
\end{theorybox}

\begin{implbox}
\begin{itemize}
  \item \pyfunc{true\_observable\_spectrum(E\_nu\_grid\_MeV, flux\_e, flux\_x,
    cross\_section\_e, cross\_section\_x, n\_target)} implements
    \labelcref{eq:detector-true-rate} by trapezoidal quadrature
    (\code{torch.trapezoid}) over the neutrino-energy axis; accepts a batched
    leading shape on \code{flux\_e}/\code{flux\_x} (e.g. a batch of
    oscillation-parameter draws), broadcasting a single cross-section grid
    against all of them.
  \item \pyfunc{apply\_response(true\_spectrum, T\_grid\_MeV,
    response\_matrix)} implements \labelcref{eq:detector-response} as a
    matrix convolution: \code{response\_matrix} is $R(T'|T)$ pre-evaluated on
    a \code{(n\_Tp, n\_T)} grid (\texttt{response.py} below), and the
    integral over $T$ is again trapezoidal.
  \item \pyfunc{bin\_counts(reco\_spectrum, Tprime\_grid\_MeV,
    bin\_edges\_MeV, exposure)} implements the integral part of
    \labelcref{eq:detector-bin-counts}: for each bin it restricts the
    reconstructed-grid points to those actually lying inside
    $[T'_i,T'_{i+1}]$ (the bin edges themselves need not be grid points) and
    trapezoidally integrates over that restricted subset, requiring at least
    2 such points per bin.
  \item \pyfunc{predicted\_counts(...)} is the top-level orchestration: calls
    \pyfunc{true\_observable\_spectrum} $\to$ \pyfunc{apply\_response} $\to$
    \pyfunc{detector.common.efficiency.apply\_efficiency}
    (\cref{sec:detector-efficiency} below, implementing
    \labelcref{eq:detector-efficiency}) $\to$ \pyfunc{bin\_counts}, then adds
    \code{background\_counts} if supplied -- exactly the four steps above, in
    order. This is the single function a concrete detector's own
    \texttt{event\_rate.py} calls with its target/response/efficiency/exposure
    already plugged in (\texttt{detector.borexino.event\_rate},
    \cref{ch:detector-borexino}).
\end{itemize}
\end{implbox}

\begin{notebox}
\pyfunc{predicted\_counts} currently assumes the reconstructed-observable
grid coincides exactly with the true-observable grid,
\code{response\_matrix.shape[0] == T\_grid\_MeV.shape[0]} (i.e.\
\code{Tprime\_grid\_MeV = T\_grid\_MeV}), raising otherwise. This is not a
physical restriction -- $T$ and $T'$ are conceptually distinct axes -- but a
present implementation simplification: a detector whose natural
reconstructed grid genuinely differs from its true-energy grid must resample
one onto the other before calling in.

Every intermediate tensor in this module carries the gradient from whichever
upstream oscillation parameters produced \code{flux\_e}/\code{flux\_x}; every
other argument (cross sections, response matrix, efficiency, target,
exposure) is treated as a grid-only constant multiplied in, so autograd
differentiates the whole chain exactly like any other product of a
differentiable factor and detached constants -- this is what lets a fit
backpropagate a $\chi^2$ built from \pyfunc{predicted\_counts} directly into
oscillation parameters.
\end{notebox}

\section{\texttt{observation.py} --- a detector-agnostic measurement container}

\modulehead{detector/common/observation.py}{\pyclass{Observation}: value,
one-sided uncertainty, and an energy/bin grid for a published measurement,
independent of which detector produced it.}

\begin{implbox}
\pyclass{Observation} formalises, into a single reusable type, the ad hoc
\code{pandas.read\_csv(...)} plus column-selection pattern that would
otherwise be repeated in every notebook comparing a prediction against
published data. It stores one label, one $x$-grid value (an energy or a
bin-centre energy), a central value, and one-sided lower/upper
uncertainties per entry, plus an optional per-entry bin width:
\begin{itemize}
  \item \pyfunc{Observation.from\_symmetric(...)}: convenience constructor
    for the common case of a single symmetric uncertainty column
    (\code{sigma\_minus is sigma\_plus}).
  \item \pyfunc{Observation.bin\_edges\_MeV} (property): for a genuinely
    binned spectrum (\code{bin\_width\_MeV} set), returns the $n+1$ bin
    edges assuming each bin is centred on its \code{x\_MeV} entry,
    \begin{equation}
      \bigl[x_i - \tfrac12 w_i,\; x_i + \tfrac12 w_i\bigr],
    \end{equation}
    with no requirement that adjacent bins be contiguous (a gap or overlap
    is the responsibility of whoever built \code{x\_MeV}/\code{bin\_width\_MeV}).
    Raises if \code{bin\_width\_MeV} is \code{None} -- a pointwise
    observation (e.g.\ four representative $P_{ee}(E)$ points) has no bin
    edges to speak of.
\end{itemize}
This same container is used both for a pointwise survival-probability
measurement and for a fully binned rate spectrum (\texttt{detector.borexino.io},
\cref{ch:detector-borexino}); it does not itself know which detector or
which observable produced the numbers it holds.
\end{implbox}

\section{\texttt{target.py} --- target particle count from stoichiometry}

\modulehead{detector/common/target.py}{\pyfunc{n\_electrons}: total target
electron count, derived from molecular composition and target mass rather
than quoted directly.}

\begin{theorybox}
$N_{\rm target}$ in \labelcref{eq:detector-true-rate} above is, for an
elastic-scattering target (\cref{ch:detector-interaction}), the total number
of electrons in the instrumented mass. Rather than quoting a detector's
published electron count directly, it is derived from basic chemistry: a
target mass $m$ built from molecules of composition
$\{(\text{element},\,\text{count})\}$ and molar mass $M$ contains
\begin{equation}
  n_{\rm molecules} = \frac{m}{M}\,N_A
\end{equation}
molecules, each contributing $Z_{\rm molecule}=\sum_{\rm element}
Z_{\rm element}\times{\rm count}_{\rm element}$ electrons (the sum of atomic
numbers over the molecule's constituent atoms, for neutral matter), giving
\begin{equation}
  \keyresult{N_e = \frac{m}{M}\,N_A\,Z_{\rm molecule}.}
\end{equation}
A stoichiometric derivation is independently checkable against textbook
atomic numbers and molar masses -- the same preference this project already
applies to \texttt{medium.solar.io}'s composition-derived neutron density
(\cref{ch:solar}) over a hardcoded ratio.
\end{theorybox}

\begin{implbox}
\begin{itemize}
  \item \pyfunc{ATOMIC\_NUMBER}: a small lookup table (electrons per neutral
    atom $=Z$) for the light elements making up common
    liquid-scintillator/water targets (H, C, N, O) -- textbook values, not
    detector-specific measurements.
  \item \pyfunc{n\_electrons(composition, molar\_mass\_g\_mol, mass\_ton)}
    implements $N_e=(m_{\rm ton}\times10^6\,{\rm g/ton}/M)\,N_A\,Z_{\rm molecule}$,
    with \code{composition} a mapping element $\to$ atom count in one target
    molecule (e.g.\ \code{\{"C": 9, "H": 12\}} for pseudocumene, C$_9$H$_{12}$,
    \cref{ch:detector-borexino}). Raises if any element is outside
    \pyfunc{ATOMIC\_NUMBER}, or if \code{molar\_mass\_g\_mol}/\code{mass\_ton}
    is non-positive.
\end{itemize}
\end{implbox}

\section{\texttt{response.py} --- Gaussian energy-response construction}
\label{sec:detector-response}

\modulehead{detector/common/response.py}{\pyfunc{gaussian\_response\_matrix}:
the generic response builder shared by any detector whose energy resolution
is well described by a (possibly energy-dependent) Gaussian smearing.}

\begin{theorybox}
For a detector whose reconstructed observable is Gaussian-distributed
around the true deposit (the common case for a photoelectron-counting
scintillator or a calorimeter with symmetric resolution), the response
function of \labelcref{eq:detector-response} is a normal density in $T'$
centred on $T$, with a possibly $T$-dependent width $\sigma(T)$:
\begin{equation}
  \keyresult{R(T'|T) = \frac{1}{\mathcal{N}(T)}\,
    \exp\!\left[-\frac{1}{2}\left(\frac{T'-T}{\sigma(T)}\right)^{\!2}\right],}
\end{equation}
where $\mathcal{N}(T)=\int dT'\,\exp[-\tfrac12((T'-T)/\sigma(T))^2]$
enforces the normalisation $\int dT'\,R(T'|T)=1$ required above. On the
infinite line $\mathcal{N}(T)=\sqrt{2\pi}\,\sigma(T)$ exactly; on a finite,
discrete reconstructed grid $\mathcal{N}(T)$ is instead computed by explicit
trapezoidal integration over that grid, so probability is conserved exactly
on the grid actually used downstream, at the (negligible, for a grid wide
enough to contain the bulk of the Gaussian) cost of a small deviation from
$\sqrt{2\pi}\sigma$ near the grid's edges.
\end{theorybox}

\begin{implbox}
\pyfunc{gaussian\_response\_matrix(T\_grid\_MeV, Tprime\_grid\_MeV,
sigma\_MeV)} builds $R(T'|T)$ on the outer-product grid
\code{(Tprime\_grid\_MeV, T\_grid\_MeV)}, column-by-column: each column
(fixed true $T$) is normalised independently by trapezoidal integration over
\code{Tprime\_grid\_MeV} (\code{torch.trapezoid}), clamped at
\code{torch.finfo(...).tiny} to avoid a division by zero for a pathologically
narrow \code{sigma\_MeV}. \code{sigma\_MeV} may be a scalar (constant
resolution) or a tensor shaped \code{(n\_T,)} (energy-dependent resolution,
one value per \code{T\_grid\_MeV} entry); the whole function runs under
\code{torch.no\_grad()} since the response matrix is a fixed instrumental
property, not a differentiable fit parameter, in every use of it in this
project so far. A detector with a Monte-Carlo migration matrix instead of an
analytic Gaussian (e.g.\ a Cherenkov detector's reconstructed-vs-true-energy
matrix) supplies its own matrix directly to
\pyfunc{event\_rate.apply\_response} and does not need this module at all.
\end{implbox}

\section{\texttt{efficiency.py} --- detection/selection efficiency}
\label{sec:detector-efficiency}

\modulehead{detector/common/efficiency.py}{\pyfunc{apply\_efficiency} and
\pyfunc{step\_efficiency}, applied after response smearing, in
reconstructed-energy space.}

\begin{theorybox}
As derived in \texttt{event\_rate.py} above, the efficiency $\varepsilon(T')$
is a multiplicative correction acting on the reconstructed spectrum,
\labelcref{eq:detector-efficiency}. The simplest possible efficiency curve
-- and the natural first model before any finer-grained cut (fiducial
volume, pulse-shape discrimination, ...) is added -- is a hard analysis
threshold, a Heaviside step in $T'$:
\begin{equation}
  \keyresult{\varepsilon(T') = \Theta(T' - T_{\rm thr})
    = \begin{cases} 0 & T' < T_{\rm thr} \\ 1 & T' \geq T_{\rm thr}. \end{cases}}
\end{equation}
\end{theorybox}

\begin{implbox}
\begin{itemize}
  \item \pyfunc{apply\_efficiency(reco\_spectrum, efficiency)} implements
    \labelcref{eq:detector-efficiency} as a plain elementwise product;
    validates that \code{efficiency} broadcasts against
    \code{reco\_spectrum}'s final axis and that every entry lies in
    $[0,1]$.
  \item \pyfunc{step\_efficiency(Tprime\_grid\_MeV, threshold\_MeV)} builds
    $\varepsilon(T')$ above directly on a grid, under
    \code{torch.no\_grad()}: as a grid-only construction (not a fit
    parameter) it never needs to be differentiated through, so the hard
    discontinuity at $T_{\rm thr}$ is not a concern here -- unlike a
    threshold entering a likelihood as a free parameter, which would need a
    smoothed (e.g.\ sigmoid) efficiency curve instead.
\end{itemize}
\end{implbox}

\section{\texttt{background.py} --- background-rate placeholder}

\modulehead{detector/common/background.py}{\pyfunc{zero\_background}: the
explicit, zero, placeholder for $N_i^{\rm bkg}$ in
\labelcref{eq:detector-bin-counts} above.}

\begin{implbox}
No real per-experiment background model is implemented anywhere in the
project yet. A real detector's background budget (radioactive contaminants
intrinsic to the target, cosmogenic isotopes, external gammas, ...) is
itself a dedicated analysis in the source experimental papers, not something
to approximate here without a citable model (see
\texttt{detector.borexino.backgrounds}, \cref{ch:detector-borexino}, for the
concrete list of contaminants relevant to that detector). \pyfunc{zero\_background(n\_bins,
device=None, dtype=torch.float64)} returns a shape-\code{(n\_bins,)} zero
tensor, ready for \pyfunc{event\_rate.predicted\_counts}'s
\code{background\_counts} argument -- explicit and zero, not silently
omitted, so that a future Poisson fit against real spectral data knows
exactly what its background term is (and is not yet) accounting for.
\end{implbox}
